\chapter{Fundamentals of Dysarthric Speech Recognition Systems}

Recent advancements in speech recognition have significantly improved interaction with machines, yet challenges remain when dealing with atypical speech patterns such as dysarthria. Understanding the underlying concepts of speech processing, neural models, and data representation is essential before developing a robust system. This chapter presents the foundational theories required for the project, including speech signal processing, deep learning models, and evaluation metrics. It also introduces the tools and frameworks used in implementation, providing the necessary background for the system design discussed in the next chapter.

\section{Speech and Dysarthria Fundamentals}

Speech is a complex acoustic signal generated through coordinated movements of the vocal tract. Dysarthria disrupts this coordination, resulting in variations in pitch, timing, articulation, and intensity. These irregularities make it difficult for conventional ASR systems to model speech accurately.

Dysarthric speech is characterized by:
\begin{itemize}
\item Slurred or imprecise articulation
\item Irregular speech rate and rhythm
\item Reduced intelligibility
\item High variability across speakers
\end{itemize}

These characteristics make dysarthric speech an out-of-distribution problem for standard ASR models.

\section{Automatic Speech Recognition (ASR)}

Automatic Speech Recognition systems convert spoken language into text. Traditional ASR systems relied on Hidden Markov Models (HMMs) and Gaussian Mixture Models (GMMs), whereas modern systems use deep learning architectures.

End-to-end ASR models directly map audio signals to text using neural networks. Whisper, a transformer-based ASR model, has demonstrated strong robustness across languages and noise conditions. However, it is primarily trained on normal speech, leading to performance degradation on dysarthric inputs.

\section{Speech Feature Extraction}

Raw audio signals are transformed into feature representations before being processed by models. Common representations include:
\begin{itemize}
\item Spectrograms
\item Mel Spectrograms
\item Log-Mel Features
\end{itemize}

Whisper uses log-Mel spectrogram features as input, enabling efficient learning of temporal and frequency patterns in speech.

\section{Data Augmentation and Synthetic Speech}

Due to limited availability of dysarthric speech datasets, data augmentation is widely used to improve model robustness. Techniques include:
\begin{itemize}
\item Noise injection
\item Pitch shifting
\item Time stretching
\item Speed perturbation
\end{itemize}

Synthetic speech generation using Text-to-Speech (TTS) models and voice conversion techniques further enhances dataset diversity. These methods simulate dysarthric characteristics and help improve generalization.

\section{Deep Learning Models and Fine-Tuning}

Transformer-based models have become the standard for speech recognition. Whisper employs an encoder-decoder transformer architecture for sequence-to-sequence learning.

Fine-tuning adapts a pre-trained model to a specific domain. Parameter-efficient techniques such as Low-Rank Adaptation (LoRA) allow fine-tuning with reduced computational cost by updating only a subset of model parameters.

\section{Evaluation Metrics}

The performance of speech recognition systems is evaluated using:
\begin{itemize}
\item Word Error Rate (WER)
\item Character Error Rate (CER)
\end{itemize}

WER is defined as the ratio of substitutions, deletions, and insertions to the total number of words. Lower WER indicates better performance.

\section{Programming Environment and Tools}

The implementation of the project uses:
\begin{itemize}
\item Python for model development
\item PyTorch for deep learning
\item Hugging Face Transformers for Whisper implementation
\item Librosa for audio processing
\item PEFT (LoRA) for efficient fine-tuning
\end{itemize}

These tools provide flexibility for experimentation and deployment.

\section{Summary}

This chapter introduced the fundamental concepts required for dysarthric speech recognition, including speech processing, ASR systems, data augmentation, and evaluation metrics. It also outlined the tools and techniques used for implementation. These foundations support the system architecture and design methodology presented in the next chapter.